% from main.tex
%%%%%%%%%%  scan idx 274  |  folio 261  %%%%%%%%%%
type \(\mathcal O_{U,X}\), où \(U\) est un objet de \(X\). Utilisant l'exactitude à droite et la commutation aux sommes directes des foncteurs \(H^p(E)\otimes\,\cdot\,\) et \(H^p(E\otimes^{\mathbf L}\,\cdot\,)\), on se ramène à prouver que la flèche de a') est un isomorphisme pour \(M=\mathcal O_{U,X}\), \(U\in\Ob X\). Or \(H^p(E)\otimes\mathcal O_{U,X}\) (resp. \(H^p(E\otimes^{\mathbf L}\mathcal O_{U,X})\)) s'identifie canoniquement (SGA 4 IV 2.13) à \(i_!(H^p(E)|U)\) (resp. \(H^p(i_!(E|U))\)), où \(i_!\) est le prolongement par \(0\) en dehors de \(U\), et l'on gagne car \(i_!\) est exact.

Supposons \(E\) pseudo-cohérent, et prouvons \(\mathrm b)\Longrightarrow\mathrm b')\). Il suffit évidemment de montrer que \((*)\) est un isomorphisme pour \(M\) injectif. Comme la flèche \((***)\) est un isomorphisme pour \(M\) injectif, l'assertion va donc résulter de la suivante :

\((5.1.1)\) : la flèche \((**)\) est un isomorphisme si \(E\) est pseudo-cohérent et \(M\) injectif.

\emph{Preuve de \((5.1.1)\).} Notons \(T\) (resp. \(T'\)) le foncteur de \(D^-(X)\) dans \(D(X)\) défini par \(F\mapsto F\otimes^{\mathbf L}M\) (resp. \(\RHom(F^{\vee},M)\)). Soit \(n\in\mathbb Z\) ; si \(H^i(F)=0\) pour \(i>n\), alors \(H^iT(F)=0\) pour \(i>n\), \(H^i(F^{\vee})=0\) pour \(i<-n\), et \(H^iT'(F)=0\) pour \(i>n\) si \(M\) est injectif. Comme d'autre part la flèche \((**)\) commute à la restriction à un ouvert, et qu'un injectif restreint à un ouvert est encore injectif, on peut, quitte à se localiser et à tronquer, supposer \(E\) strictement parfait (2.1). Par dévissage, on se ramène alors à \(E=\mathcal O_X\), et l'on gagne.

\textbf{Remarques 5.2.} (i) La condition a) (resp. a'), b), b')) est de nature locale : si \((X_\alpha)\) est un recouvrement ouvert de l'objet final de \(X\), et si, pour chaque \(\alpha\), la condition a) (resp\ldots) est vérifiée pour \(X_\alpha\) et \(E\,|\,X_\alpha\), alors a) (resp\ldots) est vérifiée.

(ii) Si \(x\) est un point de \(X\), notons \(\mathrm a_x)\) (resp. \(\mathrm a'_x)\), etc.) la
